\centerline{\bf \Huge Вокруг ортоцентра}

Основные свойства.
\begin{enumerate}
    \item Направления из вершины треугольника на ортодентр и центр ошисанной окружности симметричны относительно биссектрисы соответствующего угла треугольника;

    \item Точки, симметричные ортоцентру относительно сторон треугольника, лежат на описанной окружности данного треугольника;

    \item Точки, симметричные ортоцентру относительно середин сторон треугольника, лежат на описанной окружности данного треугольника, причем являются диаметрально шротивоположными к вершинам треугольника;

    \item Высоты треугольника являются биссектрисами ортотреугольника, то есть треугольника, вершинами которого являются основания высот исходного треугольника.
\end{enumerate}



\problem{1} Докажите, что расстояние от вершины треугольника до точки пересечения высот вдвое больше, чем расстояние от центра описанной окружности до противоположной стороны.

\problem{2} Высоты $A D$ и $C E$ остроугольного треугольника $A B C$ пересекаются в точке $H$. Окружность, описанная около треугольника $A C H$, пересекает стороны $A B$ и $B C$ в точках $F$ и $G$ соответственно. Докажите, что $F G=2 D E$.

\problem{3}
$\bigl[\textbf{Олимпиада Шарыгина, заочный тур, 2023}\bigr]$
В треугольнике $A B C$ проведены высоты $A H_A$ и $B H_B$. Прямая $H_A H_B$ пересекает описанную окружность треугольника $A B C$ в точках $P$ и $Q$. Точка $A^{\prime}$ симметрична точке $A$ относительно $B C$, точка $B^{\prime}$ симметрична точке $B$ относительно $C A$. Докажите, что $A^{\prime}, B^{\prime}, P, Q$ лежат на одной окружности.

\problem{4} 
$\bigl[\textbf{Московская устная олимпиада по геометрии, 2014}\bigr]$
Отрезок $A D$ - диаметр описанной окружности остроугольного треугольника $A B C$. Через точку $H$ пересечения высот этого треугольника провели прямую, параллельную стороне $B C$, которая пересекает стороны $A B$ и $A C$ в точках $E$ и $F$ соответственно. Докажите, что периметр треугольника $D E F$ в два раза больше стороны $B C$.

\newpage

\centerline{\bf \Huge }
\centerline{\bf \Huge }


\problem{5} Через точку пересечения высот остроугольного треугольника $A B C$ проходят три окружности, каждая из которых касается одной из сторон треугольника в основании высоты. Докажите, что вторые точки пересечения окружностей являются вершинами треугольника, подобного исходному.

\problem{6}
$\bigl[\textbf{Олимпиада Шарыгина, заочный тур, 2023}\bigr]$
Про треугольник $A B C$ известно, что точка, симметричная ортоцентру относительно центра описанной окружности, лежит на стороне $B C$. Пусть $A_1$ - основание высоты, проведенной из точки $A$. Докажите, что $A_1$ лежит на окружности, проходящей через середины трех высот треугольника $A B C$.

\problem{7} 
$\bigl[\textbf{Олимпиада Шарыгина, закл, 2007}\bigr]$
Точки $A^{\prime}, B^{\prime}, C^{\prime}$ - основания высот остроугольного треугольника $A B C$. Окружность с центром $B$ и радиусом $B B^{\prime}$ пересекает прямую $A^{\prime} C^{\prime}$ в точках $K$ и $L$ (точки $K$ и $A$ лежат по одну сторону от $B B^{\prime}$ ). Докажите, что точка пересечения прямых $A K$ и $C L$ лежит на прямой $B O$, где $O$ - центр описанной окружности треугольника $A B C$.